\chapter{Progettazione Fisica}

\section{Il Livello Fisico: Definizione dei Metadati e Linguaggio SQL}

Conclusa la fase di progettazione logica, il progettista ha tra le mani uno schema relazionale formale e matematicamente inattaccabile. Tuttavia, questo schema risiede ancora su carta. Il passo successivo consiste nell'istanziare fisicamente questo modello all'interno di un \textbf{DBMS} (Database Management System, come PostgreSQL, MySQL o Oracle).

Per comunicare con il motore del database si utilizza \textbf{SQL (Structured Query Language)}, il linguaggio standard universale per la gestione dei database relazionali.

\begin{dettagli}[Lo Standard SQL e la Portabilità]
Sebbene SQL sia il linguaggio standard per i database relazionali (codificato dagli enti internazionali ANSI e ISO), nella realtà commerciale ogni motore DBMS (PostgreSQL, MySQL, Oracle, ecc.) implementa un proprio "dialetto", introducendo estensioni e funzioni proprietarie non previste dallo standard. 

Per questo motivo, il manuale ufficiale di un buon DBMS inizia solitamente dichiarando esplicitamente il proprio livello di aderenza allo standard internazionale. Dal punto di vista dell'ingegneria del software, scrivere codice il più aderente possibile allo standard puro è una \textit{best practice} fondamentale: assicura un'elevata \textbf{portabilità} del sistema, garantendo la possibilità di migrare l'infrastruttura da un DBMS a un altro con sforzi e costi di riscrittura minimi. Nel corso della progettazione si farà riferimento esclusivamente alle direttive dell'SQL standard.
\end{dettagli}

\subsection*{Il concetto di Metadato e il Catalogo di Sistema}

Prima di poter inserire i dati reali dell'applicazione (es. Mario Rossi, matricola 12345), il DBMS deve conoscere la "forma" che questi dati avranno, le regole a cui devono sottostare e i legami che li uniscono. 
Queste informazioni strutturali prendono il nome di \textbf{Metadati} (letteralmente: \textit{dati che descrivono altri dati}).

Quando si definiscono gli schemi e i vincoli in SQL, il motore del database non sta ancora salvando informazioni applicative, ma sta popolando il proprio \textbf{Catalogo di Sistema} (o \textit{Data Dictionary}). Il catalogo è un set di tabelle interne, invisibili all'utente finale, che il DBMS usa per ricordare quante tabelle esistono, come si chiamano le colonne, quali sono i tipi di dato ammessi e quali vincoli di chiave primaria o esterna devono essere imposti.

\subsection*{L'Anatomia di SQL: Un Linguaggio Dichiarativo}

A differenza di linguaggi imperativi come Java, C o Python, in cui il programmatore deve descrivere \textit{passo dopo passo} come eseguire un'operazione, SQL è un linguaggio \textbf{dichiarativo}: il progettista si limita a dichiarare \textit{cosa} vuole ottenere o \textit{come} deve essere fatta la struttura; sarà il motore di ottimizzazione interno del DBMS a decidere la strategia algoritmica migliore per allocare la memoria o estrarre i dati.

Per gestire l'intero ciclo di vita del dato, il linguaggio SQL è suddiviso in tre grandi sotto-insiemi funzionali (sub-languages):

\begin{enumerate}
    \item \textbf{DDL (Data Definition Language):}
    È il linguaggio di definizione dei metadati. Viene utilizzato esclusivamente per manipolare la struttura del database (il contenitore). I suoi comandi principali sono:
    \begin{itemize}
        \item \texttt{CREATE}: per generare nuove tabelle, indici o viste.
        \item \texttt{ALTER}: per modificare la struttura di una tabella esistente (es. aggiungere una colonna).
        \item \texttt{DROP}: per distruggere fisicamente una struttura e i suoi dati.
    \end{itemize}

    \item \textbf{DML (Data Manipulation Language):}
    È il linguaggio utilizzato per manipolare il contenuto (i dati veri e propri) all'interno delle strutture precedentemente create col DDL. I comandi principali sono:
    \begin{itemize}
        \item \texttt{INSERT}: per inserire nuove tuple.
        \item \texttt{UPDATE}: per aggiornare i valori di tuple esistenti.
        \item \texttt{DELETE}: per rimuovere tuple.
    \end{itemize}

    \item \textbf{DQL (Data Query Language):}
    Spesso considerato parte del DML, è l'istruzione \texttt{SELECT}, il cuore pulsante dei database relazionali, utilizzata per interrogare il motore e filtrare, aggregare o unire (\texttt{JOIN}) i dati.
\end{enumerate}

\section{I Domini di Base (Tipi di Dato) nello Standard SQL}

Nella creazione di una tabella tramite DDL, ogni attributo deve essere obbligatoriamente associato a un \textbf{Dominio} (tipo di dato). Il dominio definisce lo spazio dei valori ammissibili per quella colonna, lo spazio fisico occupato su disco e le operazioni matematiche o logiche consentite.

Lo Standard SQL classifica i tipi di dato in grandi famiglie. Di seguito vengono presentati i domini fondamentali e maggiormente utilizzati.

\subsubsection*{Dati di tipo Stringa (Caratteri)}
Servono per memorizzare testo, codici alfanumerici o numeri su cui non si devono effettuare operazioni matematiche (es. numeri di telefono, CAP, codici fiscali).

\begin{itemize}
    \item \texttt{\textbf{CHAR(n)}}: Stringa a lunghezza \textbf{fissa} pari a $n$ caratteri. Se si inserisce una stringa più corta, il DBMS la "riempie" (padding) con spazi vuoti fino a raggiungere $n$.
    \begin{itemize}
        \item \textit{Uso ingegneristico:} Estremamente veloce per la ricerca. Si usa \textbf{solo} quando il dato ha lunghezza rigidamente costante (es. \texttt{CHAR(16)} per il Codice Fiscale, \texttt{CHAR(2)} per la sigla della Provincia).
    \end{itemize}
    
    \item \texttt{\textbf{VARCHAR(n)}}: Stringa a lunghezza \textbf{variabile} con limite massimo di $n$ caratteri. Se si inserisce una stringa più corta, occupa su disco solo lo spazio effettivo più un marcatore di fine stringa.
    \begin{itemize}
        \item \textit{Uso ingegneristico:} Lo standard \textit{de facto} per quasi tutto il testo (es. \texttt{VARCHAR(50)} per il Nome, \texttt{VARCHAR(255)} per l'Email). Fa risparmiare enorme spazio su disco, a costo di una microscopica latenza in più durante l'allocazione della memoria.
    \end{itemize}
\end{itemize}

\subsubsection*{Dati di tipo Numerico}
SQL distingue rigorosamente tra numeri esatti (dove 1.1 + 1.1 fa esattamente 2.2) e numeri approssimati (soggetti a virgola mobile e possibili errori di precisione).

\begin{itemize}
    \item \textbf{Numerici Esatti Interi:}
    \begin{itemize}
        \item \texttt{\textbf{INT}} o \texttt{\textbf{INTEGER}}: Numero intero standard con segno (solitamente a 32 bit, range da circa -2 a +2 miliardi). È il re indiscusso delle chiavi primarie surrogate e dei contatori.
        \item \texttt{\textbf{SMALLINT}}: Numero intero "piccolo" (solitamente a 16 bit). Utile per risparmiare RAM se il range di valori è sicuramente ridotto (es. l'anno di immatricolazione, o valori da 1 a 100).
    \end{itemize}
    
    \item \textbf{Numerici Decimali Esatti:}
    \begin{itemize}
        \item \texttt{\textbf{NUMERIC(p, s)}} o \texttt{\textbf{DECIMAL(p, s)}}: Numeri a virgola fissa. Il parametro $p$ (\textit{precision}) indica il numero totale di cifre, mentre $s$ (\textit{scale}) indica quante di queste cifre stanno dopo la virgola. 
        \item \textit{Esempio:} \texttt{NUMERIC(5, 2)} può memorizzare valori da -999.99 a 999.99.
        \item \textit{Uso ingegneristico:} Obbligatorio per i dati \textbf{finanziari e monetari} (stipendi, prezzi, bilanci). Garantisce che non ci siano mai arrotondamenti imprevisti.
    \end{itemize}
    
    \item \textbf{Numerici Approssimati:}
    \begin{itemize}
        \item \texttt{\textbf{REAL}}, \texttt{\textbf{DOUBLE PRECISION}}, \texttt{\textbf{FLOAT}}: Numeri a virgola mobile (standard IEEE 754).
        \item \textit{Uso ingegneristico:} Si usano \textbf{solo} per calcoli scientifici, coordinate GPS o misurazioni fisiche dove la scala è immensa e una micro-approssimazione (es. 1.9999999 al posto di 2) è tollerabile. Da evitare assolutamente per la valuta.
    \end{itemize}
\end{itemize}

\subsubsection*{Dati di tipo Temporale}
\begin{itemize}
    \item \texttt{\textbf{DATE}}: Memorizza solo la data (Anno, Mese, Giorno). Formato: \texttt{'YYYY-MM-GG'}. Ottimo per date di nascita o date di assunzione.
    \item \texttt{\textbf{TIME}}: Memorizza solo l'orario (Ore, Minuti, Secondi). Formato: \texttt{'HH-MI-SS'}.
    \item \texttt{\textbf{TIMESTAMP}} o \texttt{\textbf{DATETIME}}: Memorizza data e ora esatte in un unico campo. Formato: \texttt{'YYYY-MM-GG HH-MI-SS'}. Fondamentale per i log di sistema, la data di creazione di un record o per tracciare transazioni.
\end{itemize}

\begin{dettagli}[La Rappresentazione Fisica dei Tipi Temporali]
Nonostante in fase di interrogazione e inserimento le date vengano espresse in formato testuale (es. \texttt{'YYYY-MM-DD'}), a livello di immagazzinamento fisico i tipi temporali \textbf{non sono stringhe}. 

Il DBMS analizza la stringa immessa e la converte in un numero (spesso un intero a 32 o 64 bit che rappresenta i giorni o i millisecondi trascorsi da una specifica data di riferimento, detta \textit{Epoch}). 
Questa astrazione matematica è cruciale per tre motivi:
\begin{enumerate}
    \item \textbf{Performance aritmetica:} Permette al DBMS di sommare giorni o calcolare intervalli tramite velocissime operazioni matematiche a livello di CPU, evitando il costoso \textit{parsing} testuale.
    \item \textbf{Ottimizzazione spaziale:} Un \texttt{DATE} occupa tipicamente 4 byte, contro i 10 byte di un \texttt{CHAR(10)}.
    \item \textbf{Validazione semantica:} Il DBMS controlla nativamente l'esistenza reale della data (es. impedendo l'inserimento del 30 Febbraio).
\end{enumerate}
\end{dettagli}

\subsubsection*{Dati di tipo Booleano}
\begin{itemize}
    \item \texttt{\textbf{BOOLEAN}}: Ammette solo i valori di verità \texttt{TRUE} o \texttt{FALSE} (più il valore \texttt{NULL} se il campo non è obbligatorio). Ottimo per i flag (es. \texttt{is\_attivo}, \texttt{is\_pagato}).
\end{itemize}

\begin{warningbox}[Il NULL non è un tipo di dato]
È fondamentale ricordare che \texttt{NULL} non appartiene a nessun dominio specifico, ma è uno \textbf{stato} universale. \texttt{NULL} in SQL non significa "zero" né "stringa vuota", ma rappresenta "informazione mancante, sconosciuta o inapplicabile"; intuitivamente, si può pensare come una variabile incognita (come in una equazione matematica). Qualsiasi tipo di dato descritto sopra, in assenza di vincoli, può assumere lo stato \texttt{NULL}.
\end{warningbox}

\begin{dettagli}[Il supporto al tipo BOOLEAN e l'approccio legacy]
Sebbene lo standard SQL preveda nativamente il dominio \texttt{BOOLEAN}, nella pratica ingegneristica è necessario tenere conto delle limitazioni specifiche di alcuni DBMS. Il caso storico più celebre è Oracle Database, che non ha supportato il tipo booleano a livello di tabella fino alla versione 23c (2023).

Quando si opera su sistemi legacy o DBMS privi di questo dominio, la \textit{best practice} architetturale consiste nel "simulare" il booleano utilizzando un tipo di dato minimale, come \texttt{CHAR(1)} o \texttt{SMALLINT}/\texttt{NUMBER(1)}, e "blindarlo" tramite un vincolo di colonna personalizzato (il vincolo \texttt{CHECK}, che vedremo a breve) che limiti esplicitamente l'inserimento a due soli valori convenzionali.
\end{dettagli}

\section{Operatori Logici e di Confronto}

Per imporre regole di business (tramite vincoli \texttt{CHECK}) o per filtrare i dati durante le interrogazioni, SQL utilizza espressioni condizionali. Un'espressione condizionale valuta il contenuto di una cella e restituisce un risultato booleano a tre vie: \texttt{TRUE}, \texttt{FALSE} oppure \texttt{UNKNOWN} (se coinvolge un valore \texttt{NULL}).

Di seguito i principali operatori standard previsti dal linguaggio:

\subsection*{Operatori di Confronto Standard}
Si applicano a stringhe, numeri e date:
\begin{itemize}
    \item \texttt{=} (Uguale)
    \item \texttt{<>} oppure \texttt{!=} (Diverso)
    \item \texttt{<}, \texttt{>}, \texttt{<=}, \texttt{>=} (Minore, Maggiore e relativi uguali)
\end{itemize}

\subsection*{Operatori Logici (Algebra di Boole)}
Permettono di combinare più condizioni elementari:
\begin{itemize}
    \item \texttt{AND}: Restituisce \texttt{TRUE} solo se \textit{tutte} le condizioni combinate sono vere.
    \item \texttt{OR}: Restituisce \texttt{TRUE} se \textit{almeno una} delle condizioni è vera.
    \item \texttt{NOT}: Inverte il risultato logico della condizione che lo segue.
\end{itemize}

\subsection*{Operatori Speciali di SQL}
Sono costrutti sintattici avanzati che semplificano enormemente la scrittura del codice, ottimizzando le performance del motore di ricerca:
\begin{itemize}
    \item \texttt{\textbf{BETWEEN} x \textbf{AND} y}: Verifica se un valore è compreso in un intervallo continuo (inclusi gli estremi \texttt{x} e \texttt{y}). Ottimo per le date o i range numerici.
    \begin{itemize}
        \item \textit{Esempio:} \texttt{voto BETWEEN 18 AND 30} equivale a \texttt{voto >= 18 AND voto <= 30}.
    \end{itemize}
    
    \item \texttt{\textbf{IN} (val1, val2, \dots)}: Verifica se un valore appartiene a una lista di elementi discreti.
    \begin{itemize}
        \item \textit{Esempio:} \texttt{provincia IN ('RM', 'MI', 'NA')} equivale a \texttt{provincia = 'RM' OR provincia = 'MI' OR provincia = 'NA'}.
    \end{itemize}
    
    \item \texttt{\textbf{LIKE} 'pattern'}: Permette il confronto testuale parziale (pattern matching) sulle stringhe. Utilizza due caratteri jolly (wildcard):
    \begin{itemize}
        \item \texttt{\%} (Percentuale): Sostituisce una sequenza di zero o più caratteri.
        \item \texttt{\_} (Underscore): Sostituisce esattamente un singolo carattere.
        \item \textit{Esempio:} \texttt{cognome LIKE 'R\%'} (Trova tutti i cognomi che iniziano per R); \texttt{nome LIKE '\_ario'} (Trova Mario, Dario, ma non Ilario).
    \end{itemize}
\end{itemize}

\begin{warningbox}[La trappola del NULL: IS NULL vs = NULL]
Per verificare se una cella è vuota, l'intuito del programmatore suggerirebbe di scrivere \texttt{colonna = NULL}. In SQL, questa sintassi è \textbf{errata} e restituirà sempre \texttt{UNKNOWN}, facendo fallire la condizione. Poiché \texttt{NULL} rappresenta un valore "incognito", è impossibile affermare matematicamente che un'incognita sia "uguale" a un'altra incognita. 

L'unico modo legittimo in SQL per verificare l'assenza o la presenza di un dato è utilizzare gli operatori specifici di stato:
\begin{itemize}
    \item \texttt{\textbf{IS NULL}} (Verifica se il campo è vuoto)
    \item \texttt{\textbf{IS NOT NULL}} (Verifica se il campo contiene un valore)
\end{itemize}
\end{warningbox}

\section{Istruzioni DDL}

Il Data Definition Language (DDL) permette di definire la struttura fisica del database. I due comandi fondamentali per istanziare l'architettura dei dati sono \texttt{CREATE SCHEMA} e \texttt{CREATE TABLE}.

\begin{bestpractice}[Convenzione Stilistica (Case Sensitivity)]
Lo standard SQL è \textit{case-insensitive} per quanto riguarda le parole chiave del linguaggio. Tuttavia, è convenzione universale nell'ingegneria del software scrivere le keyword in \textbf{MAIUSCOLO} (es. \texttt{CREATE TABLE}, \texttt{VARCHAR}) e gli identificatori definiti dall'utente in \textbf{minuscolo} (es. \texttt{studenti}, \texttt{matricola}). Questo garantisce la massima leggibilità del codice.
\end{bestpractice}

\subsection{Creazione di uno schema logico}
\texttt{CREATE SCHEMA} crea quello che in informatica viene chiamato un \textit{namespace} o, più semplicemente, un contenitore logico (un "recinto"). Serve a raggruppare un insieme di tabelle, viste e vincoli che appartengono alla stessa applicazione, separandoli da altri progetti presenti nello stesso server database. È importante notare che una tabella non può fluttuare nel vuoto: deve sempre risiedere all'interno di uno schema. Se non si specifica uno schema personalizzato con \texttt{CREATE SCHEMA}, molti DBMS inseriscono automaticamente le tabelle in uno schema di default chiamato comunemente \texttt{public} o \texttt{dbo}.

La sintassi di base è:
\begin{lstlisting}[language=SQL]
CREATE SCHEMA <nome_schema>;
\end{lstlisting}

\subsection{Creazione di una tabella}
Il comando \texttt{CREATE TABLE} istanzia fisicamente uno schema relazionale $S(A_1, \dots, A_k)$ associando a ciascun attributo $A_i$ il proprio dominio $dom(A_i)$.
Il nome assegnato alla tabella deve essere \textbf{unico} all'interno dello schema (o del database, se non si usano schemi espliciti).

La struttura del comando prevede una coppia di parentesi tonde all'interno delle quali si susseguono istruzioni separate da virgole. 

\begin{example}[Creazione tabella senza vincoli espliciti]
Di seguito si riporta la traduzione in SQL dello schema relazionale relativo agli studenti. In questa prima fase esplorativa, ci limitiamo a definire i nomi degli attributi e i rispettivi domini (tipi di dato), omettendo temporaneamente i vincoli (come la chiave primaria).

\begin{lstlisting}[language=SQL, basicstyle=\ttfamily\small, keywordstyle=\color{blue}\bfseries]
CREATE TABLE studenti (
    matricola CHAR(9),
    cf        CHAR(16),
    nome      VARCHAR(20),
    cognome   VARCHAR(20),
    dataN     DATE
);
\end{lstlisting}

Logicamente, il corpo del comando si divide in due blocchi sequenziali:

\begin{enumerate}
    \item \textbf{Definizione degli Attributi:} In questa prima fase si elencano le colonne. Per ogni colonna si deve specificare il nome, il tipo di dato e, opzionalmente, dei vincoli specifici applicabili alla singola cella (vincoli intra-attributo).
    \item \textbf{Definizione dei Vincoli di Tabella:} In questa seconda fase (che vedremo subito a seguire), si esplicitano i vincoli complessi che coinvolgono più colonne contemporaneamente (come chiavi primarie composte o chiavi esterne).
\end{enumerate}

\medskip

La sintassi per la definizione di un singolo attributo (la riga base) è la seguente:
$$ \texttt{<nome\_attributo> <tipo\_attributo> [vincoli\_attributo]} $$


Si noti come la scelta dei domini rispecchi quanto discusso nel paragrafo precedente:
\begin{itemize}
    \item \texttt{matricola} e \texttt{cf} utilizzano \texttt{CHAR} poiché hanno una lunghezza fissa e rigidamente definita (rispettivamente 9 e 16 caratteri).
    \item \texttt{nome} e \texttt{cognome} utilizzano \texttt{VARCHAR} per ottimizzare lo spazio su disco, dato che la lunghezza reale del testo inserito è variabile.
    \item \texttt{dataN} (data di nascita) utilizza il tipo nativo \texttt{DATE} per garantire validazione e correttezza semantica.
\end{itemize}
\end{example}

\subsection{I Vincoli di Colonna (Intra-attributo o in-line)}

I vincoli di colonna sono regole di integrità applicate direttamente alla definizione di un singolo attributo, sulla stessa riga in cui viene dichiarato il tipo di dato. Servono a limitare i valori ammissibili nella singola cella, agendo come prima barriera contro la corruzione dei dati.

Lo Standard SQL prevede quattro vincoli principali a livello di colonna:

\subsubsection*{NOT NULL (Obbligatorietà)}
Indica che il campo non può mai assumere lo stato \texttt{NULL}. All'atto dell'inserimento di una nuova riga, l'utente è obbligato a fornire un valore valido per questa colonna.
\begin{itemize}
    \item \textit{Uso:} Si applica a tutti i dati vitali che definiscono l'entità (es. il cognome di una persona, il titolo di un libro).
\end{itemize}

\subsubsection*{UNIQUE (Univocità di colonna)}
Impone che tutti i valori presenti in quella colonna siano distinti tra loro. Il DBMS rifiuterà qualsiasi tentativo di inserire una riga con un valore già esistente in un'altra riga per quel campo.
\begin{itemize}
    \item \textit{Uso:} Si applica alle chiavi naturali candidate (es. il Codice Fiscale, l'Email). Nello standard SQL, a meno di specifiche implementazioni dei singoli DBMS, una colonna \texttt{UNIQUE} consente l'inserimento di molteplici valori \texttt{NULL}, poiché \texttt{NULL} rappresenta una incognita e due incognite non sono considerate uguali tra loro.
\end{itemize}

\subsubsection*{PRIMARY KEY (Chiave Primaria a livello di colonna)}
Elegge l'attributo a identificatore principale della tabella. Specificare \texttt{PRIMARY KEY} su una colonna implica contemporaneamente i vincoli \texttt{NOT NULL} e \texttt{UNIQUE}. Una tabella può avere \textbf{una sola} chiave primaria.
\begin{itemize}
    \item \textit{Uso:} Si usa quando la chiave primaria è composta da un solo attributo (es. una chiave surrogata come \texttt{id} o una chiave naturale singola come \texttt{matricola}).
\end{itemize}

\subsubsection*{DEFAULT (Valore predefinito)}
Non è un vincolo di integrità in senso stretto, ma una direttiva di assegnazione. Se durante un inserimento (\texttt{INSERT}) l'utente omette il valore per questa colonna, il DBMS vi inserisce automaticamente il valore specificato dopo la keyword \texttt{DEFAULT}.
\begin{itemize}
    \item \textit{Uso:} Ottimo per impostare stati iniziali (es. \texttt{is\_attivo BOOLEAN DEFAULT TRUE}) o timestamp di creazione (es. \texttt{data\_inserimento DATE DEFAULT CURRENT\_DATE}).
\end{itemize}

\begin{example}[Evoluzione dello schema studenti con vincoli di colonna]
Riprendiamo lo schema della tabella \texttt{studenti} e applichiamo i vincoli intra-attributo per blindarne la coerenza semantica ed evitare anomalie:

\begin{lstlisting}[language=SQL, basicstyle=\ttfamily\small, keywordstyle=\color{blue}\bfseries]
CREATE TABLE studenti (
    matricola CHAR(9) PRIMARY KEY,
    cf        CHAR(16) UNIQUE NOT NULL,
    nome      VARCHAR(20) NOT NULL,
    cognome   VARCHAR(20) NOT NULL,
    dataN     DATE
);
\end{lstlisting}

\textbf{Analisi delle scelte effettuate:}
\begin{itemize}
    \item \texttt{matricola}: È l'identificatore principale. Diventa \texttt{PRIMARY KEY}, quindi sarà univoca e obbligatoria per definizione.
    \item \texttt{cf}: Il codice fiscale deve essere univoco (\texttt{UNIQUE}) per evitare duplicati nella realtà, ma deve anche essere obbligatorio (\texttt{NOT NULL}) per ogni studente registrato.
    \item \texttt{nome} e \texttt{cognome}: Viene applicato \texttt{NOT NULL} perché un'istanza di studente priva di identità anagrafica non avrebbe alcun significato logico.
    \item \texttt{dataN}: Rimane priva di vincoli espliciti; questo significa che il sistema accetterà il valore \texttt{NULL} qualora la data di nascita non fosse momentaneamente disponibile al momento della registrazione.
\end{itemize}
\end{example}

\subsection{I Vincoli di Tabella (Out-of-line)}

Fino a questo momento abbiamo visto i vincoli applicati direttamente sulla stessa riga della definizione dell'attributo (sintassi \textit{in-line}). Tuttavia, lo standard SQL permette (e in molti casi impone) di dichiarare i vincoli in una sezione dedicata alla fine del comando \texttt{CREATE TABLE}, separata dalla definizione delle colonne. Questa formulazione prende il nome di \textbf{Vincolo di Tabella} (o \textit{out-of-line}).

La sintassi generale per dichiarare un vincolo a livello di tabella è:
$$ \texttt{CONSTRAINT <nome\_vincolo> <TIPO\_VINCOLO> (colonna1, colonna2, \dots)} $$

Sebbene la clausola \texttt{CONSTRAINT <nome>} sia opzionale (se omessa, il DBMS genera internamente un nome alfanumerico casuale, es. \texttt{SYS\_C00142}), nella pratica si raccomanda di assegnare sempre un nome esplicito e semantico (es. \texttt{pk\_studenti} o \texttt{unq\_cf}). 

Nominare i vincoli è fondamentale per due operazioni architetturali:
\begin{enumerate}
    \item \textbf{Leggibilità degli errori:} Se un'applicazione viola una regola, il database genererà un'eccezione specificando il nome del vincolo. Leggere \textit{"Violazione del vincolo unq\_cf"} permette al programmatore di capire istantaneamente che c'è un problema di codici fiscali duplicati, senza dover fare debug alla cieca.
    \item \textbf{Manipolazione a posteriori (Enable/Disable):} Quando il vincolo viene dichiarato, il DBMS lo crea e lo attiva automaticamente. Tuttavia, conoscendone il nome, l'amministratore del database può manipolarlo in futuro senza distruggere i dati. Tramite istruzioni specifiche (es. \texttt{ALTER TABLE ... DISABLE CONSTRAINT}), un vincolo può essere momentaneamente spento. Questa operazione è comunissima durante l'importazione massiva di milioni di dati (\textit{bulk insert}), dove spegnere i controlli accelera drasticamente la scrittura, per poi riattivarli (\textit{ENABLE}) a importazione finita.
\end{enumerate}

\subsubsection{Il caso delle Chiavi Primarie Composte}
La motivazione tecnica principale che rende indispensabile la sintassi a livello di tabella è la gestione dei vincoli multi-colonna. 

Nel nostro esempio precedente, la chiave primaria era formata da un solo attributo (\texttt{matricola}). Ma cosa accade se la progettazione logica produce una \textbf{chiave primaria composta} (es. una tabella che memorizza gli esami superati da uno studente in un determinato corso)?
In SQL è illegale scrivere \texttt{PRIMARY KEY} su due righe separate, perché il DBMS tenterebbe di creare \textit{due} chiavi primarie distinte per la stessa tabella. In questo scenario, \textbf{l'uso del vincolo di tabella diventa strutturalmente obbligatorio}, poiché permette di racchiudere più colonne sotto un'unica regola.

\begin{example}[Riscrittura e Chiavi Composte]
Mostriamo la sintassi applicata allo schema \texttt{studenti} e a una nuova tabella associativa \texttt{esami\_sostenuti}.

\begin{lstlisting}[language=SQL, basicstyle=\ttfamily\small, keywordstyle=\color{blue}\bfseries]
-- Esempio 1: Riscrittura con vincoli di tabella nominati
CREATE TABLE studenti (
    matricola CHAR(9),
    cf        CHAR(16) NOT NULL, -- NOT NULL resta solitamente in-line
    nome      VARCHAR(20) NOT NULL,
    cognome   VARCHAR(20) NOT NULL,
    dataN     DATE,
    
    -- Sezione Vincoli di Tabella
    CONSTRAINT pk_studenti PRIMARY KEY (matricola),
    CONSTRAINT uq_studenti_cf UNIQUE (cf)
);

-- Esempio 2: Obbligatorieta' del vincolo per PK composta
CREATE TABLE esami_sostenuti (
    studente_matr CHAR(9),
    codice_corso  CHAR(5),
    voto          INT NOT NULL,
    
    -- La PK unisce due attributi e deve essere definita qui
    CONSTRAINT pk_esami PRIMARY KEY (studente_matr, codice_corso)
);
\end{lstlisting}
\end{example}

\begin{bestpractice}[Convenzione e Best Practice]
Come regola pratica di sviluppo: il vincolo d'obbligatorietà (\texttt{NOT NULL}) si definisce \textbf{sempre e solo in-line} accanto all'attributo. Al contrario, i vincoli strutturali d'identità (\texttt{PRIMARY KEY}, \texttt{UNIQUE}) e quelli di navigabilità (\texttt{FOREIGN KEY}, che vedremo a breve) vengono quasi sempre traslati a fine tabella e nominati esplicitamente, per favorire l'ordine mentale e la manutenzione del codice.
\end{bestpractice}

% \subsection{Vincoli Semantici Personalizzati: La clausola CHECK}

% Spesso, i tipi di dato standard (\texttt{INTEGER}, \texttt{VARCHAR}, ecc.) non sono sufficienti per garantire la correttezza della \textit{business logic} dell'applicazione. Per imporre regole di validazione personalizzate (vincoli di dominio) si utilizza la clausola \textbf{\texttt{CHECK}}. Questa direttiva valuta un'espressione booleana: se la riga inserita rende l'espressione \texttt{TRUE}, il dato viene accettato; in caso contrario, la transazione viene rifiutata.

% \begin{example}[Tabella Esami: Applicazione di vincoli multipli]
% Sulla base dello schema concettuale, definiamo la tabella \texttt{esami}. L'obiettivo è tradurre in codice SQL tre regole di business fondamentali:
% \begin{enumerate}
%     \item \textbf{Vincolo Lode:} L'attributo \texttt{lode} è un falso booleano (\texttt{CHAR(1)}). Deve poter valere solo 'Y' (Yes) o 'N' (No).
%     \item \textbf{Vincolo Voto:} L'attributo \texttt{voto} deve essere rigorosamente compreso tra 0 e 30.
%     \item \textbf{Chiave Esterna:} La matricola inserita nell'esame deve corrispondere a uno studente realmente esistente nel database.
% \end{enumerate}

% Ecco la traduzione completa in linguaggio DDL (si assume la preesistenza della tabella \texttt{studenti} con primary key \texttt{matricola}):

% \begin{lstlisting}[language=SQL, basicstyle=\ttfamily\small, keywordstyle=\color{blue}\bfseries]
% CREATE TABLE esami (
%     matr  CHAR(9), -- Deve coincidere con la PK di studenti
%     corso VARCHAR(20) NOT NULL,
%     data  DATE NOT NULL,
%     voto  INTEGER NOT NULL,
%     -- in questo caso supponiamo che il DBMS non supporti il tipo boolean
%     lode  CHAR(1) DEFAULT 'N',
    
%     -- Definizione della Primary Key composta (studente, corso)
%     CONSTRAINT pk_esami PRIMARY KEY (matr, corso),
    
%     -- 1. Vincolo Lode (Simulazione Booleano)
%     CONSTRAINT chk_lode CHECK (lode IN ('Y', 'N')),
    
%     -- 2. Vincolo Voto (Range di validita')
%     CONSTRAINT chk_voto CHECK (voto >= 0 AND voto <= 30),
    
%     -- 3. Chiave Esterna (Integrita' Referenziale)
%     CONSTRAINT fk_esami_studenti FOREIGN KEY (matr) 
%         REFERENCES studenti(matricola)
% );
% \end{lstlisting}
% \end{example}

% \subsubsection{Analisi del Vincolo CHECK}
% L'operatore \texttt{CHECK} è estremamente flessibile e supporta i normali operatori logici SQL (\texttt{AND}, \texttt{OR}, \texttt{IN}, \texttt{BETWEEN}, \texttt{LIKE}). 
% Nel caso del \texttt{voto}, il vincolo poteva essere sintatticamente scritto anche utilizzando l'operatore di inclusione: \texttt{CHECK (voto BETWEEN 0 AND 30)}.

% \subsubsection{Introduzione alla Chiave Esterna (FOREIGN KEY)}
% L'ultimo vincolo inserito (\texttt{fk\_esami\_studenti}) rappresenta un salto di paradigma: non è più un vincolo "intra-relazionale" (che controlla i dati all'interno della stessa tabella), ma un vincolo \textbf{inter-relazionale}. Esso crea un ponte strutturale tra la tabella \texttt{esami} e la tabella \texttt{studenti}, impedendo l'inserimento di un esame associato a un numero di matricola inesistente.







